\section*{Author foundation and notation (publisher section 2)}
We work in a minimal variant of the Calculus of Inductive Constructions (CIC, [12, 11, 53]).
Concretely, like CIC, it has a hierarchy of predicative type universes $\mathcal T_i$, where we omit the index $i$, an impredicative universe of propositions $\mathcal P$, and inductive types in all universes.
We use inductive types $\mathbb B$ with constructors $\mathsf{true}$ and $\mathsf{false}$, $1$ with $\star$ as well as $\mathbb N$.
We use $\langle n,m\rangle$ as Cantor pairing on $\mathbb N$.
We use the notation $\forall$ and $\to$ for both propositions and types,
$\land,\lor,\exists$ for propositions where the constructors of $\lor$ are $\mathsf{left}_{\lor}$ and $\mathsf{right}_{\lor}$, and $\times,+,\Sigma$ for types with constructors $(a,b)$ for $\times,\exists$ and $\Sigma$, and $\mathsf{inl},\mathsf{inr}$ for $+$.
We also use inductive equality $x=y$ with sole constructor $\mathsf{refl}:x=x$.
Crucially, minimal CIC restricts case analysis on proofs, i.e.\ $H:P:\mathcal P$ in a $\mathcal T$ context. They are only allowed for falsity and equality, i.e.:
\[
\forall T:\mathcal T.\ \bot\to T\qquad
\forall T:X\to\mathcal T.\ T\,x\to x=y\to T\,y.
\]
For all other propositional types, induction and elimination principles are only provable in the form $T:X\to\cdots\to Y\to\mathcal P$.
Propositions a priori behave as other types, meaning in particular they are by default neither extensional nor proof-irrelevant. Both can however be consistently assumed:
\[
\mathsf{PE}:=\forall P\,Q:\mathcal P.\ (P\leftrightarrow Q)\to P=Q
\qquad
\mathsf{PI}:=\forall P:\mathcal P.\ \forall H\,H':P.\ H=H'.
\]
Further extensionality axioms one can impose concern functions (FE): $\mathsf{FE}:=\forall X\,Y.\ \forall f\,g:X\to Y.\ (\forall x.\ f\,x=g\,x)\to f=g$.
Finally, we distinguish the following variants of the excluded middle:
\[
\mathsf{LEM}:=\forall P:\mathcal P.\ P\lor\neg P\qquad
\mathsf{ELEM}:=\forall X:\mathcal T.\ \forall x\,y:X.\ x=y\lor x\ne y.
\]
For LEM, we speak of $P$ being \emph{definite} and reserve \emph{decidable} for sums $P+\neg P$ or functions $f:X\to\mathbb B$ such that $\forall x.\ P\,x\leftrightarrow f\,x=\mathsf{true}$.
In foundations with unique choice, the two are equivalent---other authors thus use the word decidable already for the propositional notion.
\section*{Coquand's proof irrelevance theorem (publisher section 5)}
In his seminal article on CIC, Coquand proved that LEM implies PI [11, p.\ 19] for Leibniz equality, exploiting Girard's paradox [26]. Using inductive equality rather than Leibniz equality, we refine his result by showing that assuming ELEM actually suffices, even when restricted to only apply to equalities of propositions. Our proof relies on Hedberg's theorem [30] and Lawvere's theorem [42], which we quickly recall here and otherwise paraphrases the reformulation of Barbanera and Berardi [3].

\paragraph{Theorem 18 (Hedberg).} Every type with definite equality has unique identity proofs.
\begin{proof}
Given $d:\forall x\,y:X.\ (x=y)\lor(x\ne y)$ we define a function $h:\forall x\,y.\ x=y\to x=y$ that for an input proof $e:x=y$ returns the proof obtained from invoking $d$, only using $e$ to refute the negative outcome $x\ne y$.
Then by construction $h$ is constant, i.e.\ $h\,x\,y\,e=h\,x\,y\,e'$ for all $e,e':x=y$.
Now in fact every function $g:\forall x\,y.\ x=y\to x=y$ has a left inverse $g^{-1}$ by setting $g^{-1}\,e:=\overline{g\,\mathsf{refl}};e$, where $\mathsf{refl}:x=x$ is the reflexivity proof for $x$, the symmetry operation is denoted by overline, and the transitivity operation is denoted by semicolon.
But then $h$ is both constant and injective, inducing that proofs $e,e':x=y$ are identical:
\[
e=h^{-1}(h\,e)=h^{-1}(h\,e')=e'.
\]
\end{proof}
Note that under FE, this proof applies to every type with double-negation stable equality.
\paragraph{Theorem 19 (Lawvere).} If there is a surjection $s:X\to(X\to Y)$, then every function $f:Y\to Y$ has a fixed point.
\begin{proof}
For $g:X\to Y$ defined by $g\,x:=f\,(s\,x\,x)$ we obtain $x$ with $s\,x=g$ by surjectivity of $s$.
Then $y:=s\,x\,x$ is a fixed point of $f$ because $f\,y=f\,(s\,x\,x)=g\,x=s\,x\,x=y$.
\end{proof}
Now our variant of Coquand's result can be stated and proved as follows.

\paragraph{Theorem 20.} If equality of propositions is definite, then PI holds.

\begin{proof}
Consider an inductive (or impredicatively encoded) proposition $B$ with two proofs $T$ and $F$.
Note that for PI it is enough to show $T=F$ as then given an arbitrary proposition $P$ with proofs $H$ and $H'$ we can define a function $f:B\to P$ by cases via $f\,T:=H$ and $f\,F:=H'$ for which we then observe $H=f\,T=f\,F=H'$.
We define $U:\mathcal P:=\forall P:\mathcal P.\ P\to B$ together with $s:U\to(U\to B)$ by $s\,u:=u\,U$.
This function is a surjection, as we can define its inverse $i:(U\to B)\to U$ by cases
\[
i\,H\,P:=\begin{cases}
H^e&\text{if }e:U=P\text{ and }H^e\text{ is the transport of }H\text{ along }e,\\
\lambda p.F&\text{otherwise (in fact any element of }P\to B\text{ works)}.
\end{cases}
\]
and then show that $s\,(i\,H)=i\,H\,U=H^e=H$, where in the last step we use Hedberg's theorem to argue that $e:U=U$ is $\mathsf{refl}$ and therefore transport along $e$ is the identity.
But now Lawvere's theorem applied to the surjection $s:U\to(U\to B)$ yields a fixed point for every function $B\to B$, so in particular there is one for the negation defined by $N\,T:=F$ and $N\,F:=T$, from which $T=F$ follows as desired.
\end{proof}
Note that in fact only the equality of propositions $U$ and $P$ needs to be determined to define $i$ in the above proof, meaning that already a type conditional in the sense of Harper and Mitchell [29] is enough to enable the construction.

